\scanpage{245}
\mspage{119}
% archive sheet number: 244

Les axiomes qui ont servi pour 8.4. $1^\circ,3^\circ$ \deleted{\unsure{et 2°}}\inserted{et 8.5.b)} sont
\[
 L(1\,\mathrm{bis},2,3)\unsure{'.|/.|!}
\]

Ceux qui sont nécessaires pour 8.5. \unsure{a)} \deleted{\unsure{et}} \deleted{8.5} sont \deleted{$\unsure{L(1)}$}
\[
 L(1\,\mathrm{bis},2,3,4).
\]

Enfin, ceux nécessaires pour 8.4. $2^\circ$ sont
\[
 L(1\,\mathrm{bis},2,3,4,\deleted{\unsure{6}}7)
\]
et \unsure{ils} impliquent la \uline{saturation forte de $W$} --- (cf \S{} suivant).

$\backslash$

Il me faut réparer un oubli, ayant traité \deleted{par} par par le mépris (!) les questions de composantes connexes, jusqu'à présent.

\uline{Proposition 9.1}. Supposons que $W$ satisfait Loc $(1,\inserted{2,3,}4)$ sans plus. \deleted{(axiome de saturation, axiome du carré cocartésien).} Soient $f_i:X_i\unsure{-|\to}Y_i$ $(i\in I)$ une famille \uline{finie} de flèches de $\mathrm{Cat}$. Pour que $f=\coprod f_i:\coprod X_i\to\coprod Y_i$ soit \unsure{dans} $W$, il \deleted{\unsure{faut}} s. que $\forall i\in I$, on ait $f_i\in W$. Si $W$ satisfait de plus : Loc 6 (axiome des limites inductives), alors cette conclusion reste valable sans condition de finitude sur $I$.

\uline{Dém}. Le cas $I$ fini se traite par réc. \unsure{$n$|$u$}$=\mathrm{card}\,I$, on est ramené aussitôt au cas où $n=2$. * On voit d'abord \deleted{\unsure{sur}}\inserted{de}

% slanted
\margin{Le cas infini \unsure{s'y ramène via|s'y ramènera,} l'hyp. Loc (6), \unsure{qui en est la limite|puisque c'est une limite}.}
